
All major published papers on reinforcement learning with execution feedback (RLEF) for code large language models (LLMs) use binary pass/fail reward. Under Group Relative Policy Optimization (GRPO), this choice produces zero gradient contribution in the dominant majority of early-training groups when the model cannot produce complete solutions. This paper analyzes the group-level advantage computation in GRPO and establishes that ratio reward (k/n test cases passing) mathematically guarantees non-zero advantage variance whenever any completion in a group achieves partial credit --- precisely the regime where binary reward guarantees zero variance. For a group of 8 completions with test-pass counts [0, 0, 1, 2, 0, 3, 0, 0], binary reward yields advantage variance 0.0000 while ratio reward yields 0.0475. A 1,000-group simulation under realistic early-training conditions (p\_pass = 0.1 per test case, G = 8, n = 5) shows that 987/1,000 groups (98.7\%) fall into this dead zone under binary reward. A controlled 208-step GRPO experiment using DeepSeek-Coder-6.7B-instruct on the APPS dataset with max\_new\_tokens = 512 produced reward\_mean = 0.0 and clipped\_ratio = 1.0 throughout for both binary and ratio conditions: all completions were truncated at the 512-token limit before syntactic completion, making both reward functions mathematically equivalent (both = 0 for all completions). This null result identifies a necessary experimental prerequisite --- non-zero training-set partial solve rate --- that must be satisfied before policy-level effects of reward granularity can be observed. A validated GRPO training framework (35/35 unit tests passing; 9 engineering issues documented and resolved) is released to support follow-up experiments under correctly-powered conditions.

